% Catatan: Pastikan paket berikut tersedia di main.tex:
%   \usepackage{multirow}    % untuk sel bertingkat pada tabel
%   \usepackage{longtable}   % untuk tabel panjang (tabel jenis query)
%   \usepackage{array}       % untuk p{} column type
%   \usepackage{booktabs}    % opsional untuk garis tabel
%
% Catatan Penomoran: Pada dokumen sumber PDF, terdapat tiga subbab
% yang sama-sama diberi nomor 3.2.2. Pada file ini, penomoran
% disesuaikan menjadi 3.2.2, 3.2.3, dan 3.2.4 oleh LaTeX secara otomatis.

%%\documentclass[../../main.tex]{subfiles} %% Uncomment for standalone compile

%%\begin{document} %% Uncomment for standalone compile

\chapter{Metode Penelitian}

Bab ini menjelaskan metodologi penelitian yang digunakan dalam pengembangan dan evaluasi sistem \textit{information retrieval} berbasis \textit{hybrid retrieval} dan \textit{cross-encoder reranking}. Pendekatan \textit{hybrid retrieval} menggabungkan metode \textit{dense} dan \textit{sparse retrieval}, kemudian dikombinasikan dengan mekanisme \textit{reranking} menggunakan \textit{cross-encoder} melalui skema pembobotan (\textit{alpha}) untuk mengoptimalkan relevansi hasil pencarian. Pembahasan dalam bab ini mencakup perancangan basis pengetahuan berbasis dokumen ISO/IEC 27001:2022, penyusunan \textit{query} dari dokumen kebijakan Direktorat Teknologi Informasi (DTI) UGM untuk kebutuhan evaluasi, implementasi metode \textit{retrieval} dan \textit{reranking}, serta perancangan skenario pengujian berdasarkan variasi parameter penelitian.

\section{Alat dan Bahan}

\subsection{Alat Tugas Akhir}

Alat-alat yang digunakan pada penelitian ini berupa perangkat keras maupun perangkat lunak sebagai sarana pendukung. Perangkat keras dan lunak yang digunakan berfungsi untuk menulis kode, pengembangan sistem, melakukan pengujian, dan pengumpulan data.

\begin{enumerate}
	\item Laptop Lenovo Ideapad Flex 5i dengan spesifikasi
	\begin{itemize}
		\item CPU: 12th Gen Intel(R) Core(TM) i7-1255U (10 cores)
		\item RAM: 16 GB (2x8GB DDR4 4267 MHz)
		\item GPU: Intel(R) Iris(R) Xe Graphics (128 MB VRAM)
		\item Penyimpanan: SSD 512 GB (INTEL SSDPEKNW512GZL)
		\item OS: Windows 11 Home Single Language (Build 22631)
	\end{itemize}

	\item Model AI
	\begin{itemize}
		\item \textit{Embedding model}: paraphrase-multilingual-MiniLM-L12-v2 (384 \textit{dimensions}, \textit{multilingual} termasuk Bahasa Indonesia)
		\item \textit{Reranking model}: BAAI/bge-reranker-v2-m3 (\textit{cross-encoder multilingual})
	\end{itemize}

	\item Perangkat Lunak
	\begin{itemize}
		\item Visual Studio Code sebagai \textit{Integrated Development Environment} (IDE) yang digunakan untuk menulis kode program, melakukan \textit{debugging}, serta mengelola proses pengembangan sistem.
		\item Python 3.11.5 sebagai bahasa pemrograman utama yang digunakan untuk mengimplementasikan seluruh komponen sistem.
		\item Git sebagai sistem kontrol versi yang digunakan untuk mengelola \textit{source code} serta melakukan pelacakan perubahan selama pengembangan.
	\end{itemize}
\end{enumerate}

\subsection{Bahan Tugas Akhir}

Dalam proses pengembangan dan evaluasi sistem \textit{Retrieval-Augmented Generation} (RAG), diperlukan beberapa bahan sebagai berikut:

\begin{enumerate}
	\item Dokumen standar ISO/IEC 27001:2022 ``INTERNATIONAL STANDARD ISO/IEC 27001 Third Edition 2022-10'' sebagai sumber utama pengetahuan yang digunakan dalam proses \textit{retrieval} dan pembentukan basis dokumen.
	\item Dokumen kebijakan dari Direktorat Teknologi Informasi (DTI) UGM berbentuk surat edaran dan surat keputusan Rektor Universitas Gadjah Mada yang digunakan sebagai sumber penyusunan \textit{query} dan \textit{ground truth} untuk evaluasi sistem.
\end{enumerate}

\section{Metode yang Digunakan}

\subsection{Desain Penelitian}

Desain penelitian disusun dengan menentukan variabel independen, dependen, dan kontrol agar proses pengujian dapat dilakukan secara sistematis dan terukur. Variabel independen mencakup tiga parameter utama yang dievaluasi secara berurutan: proporsi bobot \textit{scoring} antara \textit{dense} dan \textit{sparse}, jumlah kandidat pada tahap \textit{reranker}, serta proporsi bobot \textit{scoring} antara \textit{hybrid} dan \textit{reranker}. Setiap parameter pada tahap berikutnya menggunakan konfigurasi optimal dari tahap sebelumnya sehingga pencarian parameter bersifat sekuensial, bukan kombinatorik. Variabel dependen digunakan untuk mengukur performa sistem berdasarkan akurasi \textit{retrieval} dan waktu latensi. Sementara itu, variabel kontrol dijaga tetap konstan selama eksperimen untuk memastikan hasil pengujian lebih konsisten dan dapat dibandingkan secara objektif.



\begin{table}[H]
\centering
\small
\caption{Variabel Penelitian (Independen)}
\label{tab:variabel_penelitian}
\begin{tabular}{|p{1.8cm}|p{3.0cm}|p{4.6cm}|p{3.4cm}|}
\hline
\textbf{Jenis} & \textbf{Nama Variabel} & \textbf{Deskripsi} & \textbf{Indikator / Parameter} \\
\hline
Independen & Proporsi bobot \textit{scoring} \textit{dense}:\textit{sparse} & Rasio pembobotan skor \textit{dense retrieval} dan \textit{sparse retrieval} pada \textit{hybrid retrieval} & \begin{tabular}[t]{@{}l@{}}alpha\_dense $\in$\\$\{0{,}0;\ 0{,}25;$\\$0{,}5;\ 0{,}75;\ 1{,}0\}$\end{tabular} \\
\hline
Independen & Jumlah kandidat & Banyaknya dokumen kandidat yang diproses oleh \textit{cross-encoder reranker} & $k \in \{5;\ 10;\ 15;\ 20\}$ \\
\hline
Independen & Proporsi bobot \textit{scoring} \textit{hybrid}:\textit{reranker} & Rasio pembobotan skor \textit{hybrid retrieval} dan skor \textit{reranker} pada skor akhir & \begin{tabular}[t]{@{}l@{}}alpha\_reranker $\in$\\$\{0{,}0;\ 0{,}25;$\\$0{,}5;\ 0{,}75;\ 1{,}0\}$\end{tabular} \\
\hline
\end{tabular}
\end{table}

\begin{table}[H]
\centering
\small
\caption{Variabel Penelitian (Dependen dan Kontrol)}
\label{tab:variabel_penelitian_lanjutan}
\begin{tabular}{|p{1.8cm}|p{3.0cm}|p{4.6cm}|p{3.4cm}|}
\hline
\textbf{Jenis} & \textbf{Nama Variabel} & \textbf{Deskripsi} & \textbf{Indikator / Parameter} \\
\hline
Dependen & Akurasi \textit{Retrieval} & Tingkat relevansi hasil \textit{retrieval} terhadap \textit{ground truth} & Hit@3, Recall@3, MAP@3, NDCG@3 \\
\hline
Dependen & Latensi Sistem & Waktu yang dibutuhkan sistem dalam memproses hingga menghasilkan \textit{output} & Waktu eksekusi \\
\hline
Kontrol & Model \textit{Embedding} dan \textit{Reranker} & Model yang digunakan dijaga tetap sama untuk semua skenario & paraphrase-multilingual-MiniLM-L12-v2 dan BGE \textit{reranker} \\
\hline
Kontrol & \textit{Dataset} Penelitian & Kumpulan \textit{query} dan dokumen sumber yang digunakan pada seluruh skenario pengujian tetap sama; \textit{query} dibentuk dari dokumen DTI UGM & \textit{Query} DTI UGM beserta \textit{ground truth}, Dokumen ISO 27001:2022 \\
\hline
Kontrol & Lingkungan Sistem & Spesifikasi perangkat dan \textit{environment} eksekusi dijaga konstan & Laptop, OS, \textit{dependency} \\
\hline
\end{tabular}
\end{table}

\subsection{Pembangunan Basis Pengetahuan ISO/IEC 27001:2022 A.5--A.8}

Pembangunan basis pengetahuan pada penelitian ini dilakukan dengan mengumpulkan dan menyusun informasi terkait kontrol keamanan dari standar ISO/IEC 27001:2022, khususnya pada Annex A domain A.5 (\textit{Organizational Controls}), A.6 (\textit{People Controls}), A.7 (\textit{Physical Controls}), dan A.8 (\textit{Technological Controls}). Sumber utama yang digunakan adalah dokumen resmi ``INTERNATIONAL STANDARD ISO/IEC 27001 Third Edition 2022-10'', yang kemudian diperkaya dengan referensi pendukung untuk memperjelas konteks implementasi dan kebutuhan audit. Proses ini bertujuan untuk menghasilkan representasi pengetahuan yang terstruktur sehingga dapat digunakan secara optimal dalam sistem \textit{Retrieval-Augmented Generation} (RAG).

Seluruh dokumen yang digunakan sebagai sumber pengetahuan (korpus) disimpan dalam format JSON (\textit{JavaScript Object Notation}). Setiap kontrol direpresentasikan sebagai satu entitas data yang terdiri dari atribut deskriptif, indikator audit, serta kata kunci dalam dua bahasa (Bahasa Inggris dan Bahasa Indonesia). Struktur ini dirancang untuk meningkatkan kualitas pencarian semantik serta mendukung variasi \textit{query} pengguna. Struktur data JSON yang digunakan ditunjukkan pada Tabel berikut.




\begin{table}[H]
\centering
\footnotesize
\caption{Struktur Data JSON Basis Pengetahuan}
\label{tab:json_struktur_deskriptif}
\begin{tabular}{|p{3.2cm}|p{4.8cm}|p{4.8cm}|}
\hline
\textbf{Nama Atribut} & \textbf{Penjelasan} & \textbf{Contoh Isian} \\
\hline
\texttt{control\_id} & Kode unik kontrol pada ISO 27001 Annex A & A.8.32 \\
\hline
\texttt{title} & Judul dari kontrol keamanan & Manajemen perubahan \\
\hline
\texttt{objective} & Tujuan utama dari kontrol & Memastikan perubahan fasilitas pemrosesan informasi dikelola dengan aman \\
\hline
\texttt{description} & Deskripsi lengkap mengenai kontrol dan implementasinya & Perubahan sistem informasi harus tunduk pada prosedur manajemen perubahan \\
\hline
\texttt{implementation}\newline\texttt{\_guidance} & Panduan implementasi kontrol secara ringkas & Proses manajemen perubahan formal dengan penilaian dampak keamanan \\
\hline
\texttt{keywords\_en} & Kata kunci dalam Bahasa Inggris untuk mendukung \textit{retrieval} & ``change management'', ``change control'' \\
\hline
\texttt{keywords\_id} & Kata kunci dalam Bahasa Indonesia & ``manajemen perubahan'', ``kontrol perubahan'' \\
\hline
\texttt{audit}\newline\texttt{\_indicators\_en} & Indikator audit dalam Bahasa Inggris & ``unapproved changes'', ``no change records'' \\
\hline
\texttt{audit}\newline\texttt{\_indicators\_id} & Indikator audit dalam Bahasa Indonesia & ``perubahan tidak disetujui'', ``tidak ada catatan perubahan'' \\
\hline
\texttt{related}\newline\texttt{\_assets\_en} & Aset terkait dalam Bahasa Inggris & ``production systems'', ``test environment'' \\
\hline
\texttt{related}\newline\texttt{\_assets\_id} & Aset terkait dalam Bahasa Indonesia & ``sistem produksi'', ``lingkungan pengujian'' \\
\hline
\texttt{security}\newline\texttt{\_principles\_en} & Prinsip keamanan terkait dalam Bahasa Inggris & ``integrity'', ``availability'' \\
\hline
\texttt{security}\newline\texttt{\_principles\_id} & Prinsip keamanan dalam Bahasa Indonesia & ``integritas'', ``ketersediaan'' \\
\hline
\end{tabular}
\end{table}

% Catatan: Subbab berikut diberi label 3.2.2 pada dokumen sumber PDF (duplikat).
% LaTeX menomorinya secara otomatis sebagai 3.2.3.
\subsection{Pembuatan \textit{Query} dan \textit{Ground Truth}}

\textit{Query} dan \textit{ground truth} disusun sebagai komponen utama dalam proses evaluasi \textit{retrieval} sistem \textit{Information Retrieval}. \textit{Query} berperan sebagai \textit{input} yang merepresentasikan kebutuhan informasi pengguna, sedangkan \textit{ground truth} digunakan sebagai acuan kebenaran untuk menilai relevansi hasil \textit{retrieval} yang dihasilkan sistem. Dengan demikian, kombinasi keduanya memungkinkan pengukuran performa sistem secara objektif melalui perbandingan antara hasil yang dihasilkan model dan referensi yang telah ditentukan.

\textit{Query} yang digunakan dalam penelitian ini disusun berdasarkan dokumen kebijakan dan prosedur milik Direktorat Teknologi Informasi (DTI) UGM yang berkaitan dengan implementasi tata kelola dan teknologi informasi terkini. Dokumen-dokumen tersebut merepresentasikan kondisi dan kebijakan terbaru organisasi sehingga dapat menjadi acuan dalam menggambarkan status implementasi keamanan informasi UGM saat ini. Proses pembentukan \textit{query} dilakukan melalui pembacaan manual terhadap isi dokumen untuk mengidentifikasi pernyataan, kebijakan, atau deskripsi kondisi organisasi yang berkaitan dengan kontrol pada Annex A.5 hingga A.8. Pernyataan-pernyataan tersebut kemudian dicatat menjadi \textit{query} yang merepresentasikan kebutuhan pencarian informasi dalam konteks audit dan evaluasi kepatuhan keamanan informasi, beserta \textit{ground truth} yang menunjukkan kontrol ISO/IEC 27001:2022 yang relevan. Pendekatan ini dilakukan agar \textit{query} yang dihasilkan dapat merefleksikan penggunaan sistem pada kondisi dokumen organisasi yang nyata. Berikut merupakan dokumen-dokumen yang digunakan dalam penyusunan \textit{query} pada penelitian ini.

\begin{table}[H]
\centering
\caption{Dokumen DTI UGM yang Digunakan sebagai Sumber \textit{Query} Asli}
\label{tab:dti_dokumen}
\begin{tabular}{|p{7.5cm}|p{5.5cm}|}
\hline
\textbf{Nama Dokumen} & \textbf{Tentang} \\
\hline
Keputusan Rektor Universitas Gadjah Mada Nomor 357/UN1.P5/KPT/HUKOR/2026 &
Pedoman Pengelolaan Domain Universitas Gadjah Mada \\
\hline
Peraturan Rektor Universitas Gadjah Mada Nomor 13 Tahun 2025 &
Tata Kelola Teknologi Informasi di Universitas Gadjah Mada \\
\hline
Keputusan Rektor Universitas Gadjah Mada Nomor 1198/UN1.P5/KPT/HUKOR/2025 &
Standar Parameter \textit{Password} dan Keamanan Sistem Finance SIMASTER Universitas Gadjah Mada \\
\hline
Surat Edaran Rektor Universitas Gadjah Mada Nomor 3779/UN1.P/DTI/TI.02.03/2025 &
Penggunaan Perangkat Lunak di Lingkungan Universitas Gadjah Mada \\
\hline
\end{tabular}
\end{table}

% Catatan: Subbab berikut diberi label 3.2.2 pada dokumen sumber PDF (duplikat).
% LaTeX menomorinya secara otomatis sebagai 3.2.4.
\subsection{Perancangan Komponen \textit{Retrieval} dan \textit{Rerank} pada Sistem}

Perancangan komponen \textit{retrieval} pada sistem \textit{Retrieval-Augmented Generation} dilakukan untuk memastikan proses penelusuran dokumen mampu menghasilkan kandidat yang relevan secara optimal. Pendekatan utama yang digunakan adalah metode \textit{hybrid retrieval}, yaitu integrasi antara teknik \textit{sparse} dan \textit{dense} untuk memperoleh keseimbangan antara pencocokan berbasis kata kunci dan pemahaman semantik. Hasil dari tahap \textit{retrieval} awal kemudian diseleksi menjadi sejumlah kandidat teratas (top@$k$) sehingga hanya dokumen yang paling relevan yang dipertimbangkan dalam evaluasi lebih lanjut.

Untuk meningkatkan akurasi dan kualitas peringkat hasil \textit{retrieval}, digunakan pendekatan \textit{reranking} yang melakukan evaluasi ulang terhadap kandidat dokumen yang telah diperoleh. Mekanisme ini memanfaatkan model yang lebih sensitif terhadap hubungan kontekstual antara \textit{query} dan dokumen, sehingga mampu menghasilkan urutan yang lebih representatif terhadap relevansi sebenarnya. Proses \textit{reranking} juga melibatkan pengaturan konfigurasi tertentu, termasuk strategi penggabungan bobot skor (\textit{score fusion}), guna memperoleh hasil akhir yang lebih akurat, konsisten, dan sesuai dengan kebutuhan evaluasi sistem.

\subsubsection{Perancangan Komponen \textit{Retrieval} pada Sistem}

Komponen \textit{hybrid retrieval} pada penelitian ini menggunakan dua metode yang dipilih berdasarkan analisis perbandingan pada Subbab~2.3.2 dan Subbab~2.3.3.

\vspace{0.5em}
\noindent\textbf{Pemilihan Metode \textit{Sparse Retrieval}}

Berdasarkan analisis perbandingan metode \textit{sparse retrieval} pada Subbab~2.3.2, BM25 dipilih sebagai komponen \textit{sparse retrieval}. BM25 mengungguli alternatif lainnya karena memiliki model probabilistik yang kuat untuk pencocokan istilah eksplisit, normalisasi panjang dokumen bawaan, serta bukti empiris yang menunjukkan superioritasnya pada domain regulasi. Karakteristik terminologi yang eksplisit dan konsisten pada kontrol ISO/IEC 27001:2022, misalnya \textit{access control}, \textit{incident management}, dan \textit{logging}, memungkinkan pencocokan leksikal bekerja secara efektif.

\vspace{1.5em}
\noindent\textbf{Pemilihan Model \textit{Dense Embedding}}

Berdasarkan analisis perbandingan model \textit{dense embedding} pada Subbab~2.3.3, paraphrase-multilingual-MiniLM-L12-v2 dipilih sebagai komponen \textit{dense embedding}. Model ini dipilih karena menawarkan keseimbangan optimal antara kualitas representasi semantik, dukungan multilingual untuk \textit{query} berbahasa Indonesia (50+ bahasa), efisiensi komputasi dengan dimensi 384, serta bukti empiris yang menunjukkan konsistensi peningkatan kualitas \textit{dense retrieval} pada arsitektur \textit{hybrid retrieval}.

Komponen \textit{hybrid retrieval} dibangun untuk menggabungkan pencarian semantik berbasis \textit{embedding} dan pencarian leksikal berbasis BM25. Dalam rancangan Bab 3, proses ini dapat dijelaskan melalui tiga tahap utama yang sesuai dengan diagram, yaitu Basis Pengetahuan, Pemrosesan \textit{Query}, dan Pemrosesan \textit{Scoring}. Ketiga tahap tersebut menunjukkan alur dari pembentukan representasi dokumen, pemrosesan masukan pengguna, hingga penggabungan skor untuk memilih dokumen paling relevan.


\begin{figure}[H]
	\centering
	\includegraphics[width=\textwidth,keepaspectratio]{Gambar/3.6HybridComponentNew.png}
	\caption{Komponen hybrid retrieval pada sistem}
	\label{fig:retrieval_komponen}
\end{figure}

Berdasarkan alur yang ditunjukkan pada Gambar~\ref{fig:retrieval_komponen}, proses \textit{retrieval} dapat diuraikan ke dalam tiga tahap utama sebagai berikut.

\begin{enumerate}
\item \textbf{Basis Pengetahuan}

Tahap Basis Pengetahuan merupakan fase inisialisasi yang dijalankan saat sistem dimulai untuk menyiapkan seluruh representasi dokumen kontrol ISO. Pada tahap ini, data kontrol ISO dimuat dari berkas JSON, dikonversi menjadi representasi teks, kemudian diproses melalui jalur \textit{dense} untuk membentuk \textit{embedding} dokumen dan jalur \textit{sparse} untuk membangun indeks BM25. Hasil dari fase ini disimpan dalam memori agar dapat digunakan kembali ketika \textit{query} pengguna diproses.

\begin{enumerate}[label=\alph*., leftmargin=1.1cm]
\item \textbf{\texttt{\_load\_controls(controls\_path)}}: Proses ini memuat data basis pengetahuan dari berkas JSON yang berisi 93 kontrol ISO. Hasil pemuatan data menjadi daftar kontrol yang akan digunakan sebagai korpus dokumen pada proses \textit{retrieval}.

\item \textbf{\texttt{control\_to\_text(control=control)}}: Proses ini mengubah satu data kontrol ISO berbentuk JSON menjadi representasi teks datar (\textit{flat text}). Representasi teks ini menggabungkan atribut penting seperti judul, objektif, deskripsi, serta atribut pengayaan yang tersedia agar dokumen memiliki konteks yang lebih lengkap.

\item \textbf{\texttt{encode\_texts(texts=control\_texts)}}: Proses ini mengubah seluruh representasi teks kontrol ISO menjadi vektor \textit{embedding}. Vektor ini merepresentasikan makna semantik dokumen sehingga dapat dibandingkan dengan \textit{query} pada ruang vektor yang sama.

\item \textbf{\texttt{normalize(vectors=doc\_embeddings)}}: Proses ini melakukan normalisasi L2 pada vektor \textit{embedding} dokumen. Normalisasi dilakukan agar perhitungan \textit{dot product} antara \textit{query} dan dokumen setara dengan \textit{cosine similarity}.

\item \textbf{\texttt{HybridRetriever.doc\_embeddings}}: Proses ini menyimpan vektor \textit{embedding} dokumen pada atribut \texttt{doc\_embeddings}. Penyimpanan dilakukan secara \textit{in-memory} sehingga indeks vektor dapat digunakan langsung saat sistem menghitung \textit{dense score}.

\item \textbf{\texttt{tokenize(text=control\_text)}}: Proses ini mengubah representasi teks dokumen menjadi daftar token leksikal. Tokenisasi dilakukan dengan mengubah teks menjadi huruf kecil dan mengambil karakter alfanumerik agar dokumen siap diproses oleh BM25.

\item \textbf{\texttt{BM25Retriever(documents=control\_texts)}}: Proses ini menginisialisasi indeks BM25 berdasarkan seluruh teks kontrol ISO. Inisialisasi ini membentuk struktur data yang diperlukan untuk menghitung kecocokan kata kunci antara \textit{query} dan dokumen.

\item \textbf{\texttt{\_compute\_doc\_freq(tokenized\_docs)}}: Proses ini menghitung jumlah dokumen yang mengandung setiap \textit{term}. Nilai \textit{document frequency} digunakan sebagai dasar untuk menentukan seberapa umum atau spesifik suatu \textit{term} dalam korpus.

\item \textbf{\texttt{\_compute\_idf(doc\_freq)}}: Proses ini menghitung nilai \textit{inverse document frequency} untuk setiap \textit{term}. Nilai IDF memberi bobot lebih tinggi pada \textit{term} yang lebih jarang muncul sehingga dokumen dengan kata kunci spesifik dapat diberi skor lebih relevan.

\item \textbf{\texttt{BM25 index}}: Proses ini menyimpan informasi \textit{term frequency}, panjang dokumen, rata-rata panjang dokumen, dan IDF. Informasi tersebut digunakan kembali pada tahap Pemrosesan \textit{Query} untuk menghitung \textit{sparse score}.
\end{enumerate}

\item \textbf{Pemrosesan \textit{Query}}

Tahap Pemrosesan \textit{Query} dijalankan setiap kali pengguna memasukkan pertanyaan atau kebutuhan informasi. Pada tahap ini, \textit{query} diproses melalui dua jalur yang sama dengan dokumen, yaitu jalur \textit{dense} untuk menghasilkan skor kemiripan semantik dan jalur \textit{sparse} untuk menghasilkan skor kecocokan kata kunci. Hasil dari tahap ini berupa \textit{dense score} dan \textit{sparse score} untuk seluruh kontrol ISO.

\begin{enumerate}[label=\alph*., leftmargin=1.1cm]
\item \textbf{\texttt{encode\_texts(texts=[query])}}: Proses ini mengubah \textit{query} pengguna menjadi vektor \textit{embedding}. Karena fungsi yang sama digunakan pada dokumen dan \textit{query}, keduanya berada pada ruang representasi semantik yang sama.

\item \textbf{\texttt{normalize(vectors=query\_embedding)}}: Proses ini menormalkan vektor \textit{query} menggunakan normalisasi L2. Normalisasi membuat perhitungan kemiripan antara \textit{query} dan dokumen dapat dilakukan secara konsisten menggunakan \textit{dot product}.

\item \textbf{\texttt{np.dot(doc\_embeddings, q\_emb)}}: Proses ini menghitung hasil \textit{dot product} antara vektor \textit{query} dan seluruh \textit{embedding} dokumen. Hasilnya adalah \textit{dense score} yang merepresentasikan tingkat kemiripan semantik \textit{query} terhadap 93 kontrol ISO.

\item \textbf{\texttt{tokenize(text=query)}}: Proses ini mengubah \textit{query} menjadi daftar token leksikal. Tokenisasi \textit{query} menggunakan aturan yang sama dengan dokumen agar pencocokan kata kunci pada BM25 berjalan konsisten.

\item \textbf{\texttt{score\_query(query=query)}}: Proses ini menghitung skor BM25 antara \textit{query} dan seluruh dokumen. Skor dihitung berdasarkan kombinasi \textit{term frequency}, IDF, dan normalisasi panjang dokumen sehingga menghasilkan \textit{sparse score} untuk 93 kontrol ISO.
\end{enumerate}

\item \textbf{Pemrosesan \textit{Scoring}}

Tahap Pemrosesan \textit{Scoring} berfungsi untuk menggabungkan \textit{dense score} dan \textit{sparse score} menjadi satu skor akhir. Kedua jenis skor terlebih dahulu dinormalisasi karena memiliki skala yang berbeda. Setelah itu, skor digabungkan menggunakan \textit{weighted sum}, diurutkan, dan digunakan untuk memilih dokumen teratas sebagai hasil \textit{retrieval}.

\begin{enumerate}[label=\alph*., leftmargin=1.1cm]
\item \textbf{\texttt{\_minmax(scores=dense\_scores)}}: Proses ini menormalisasi \textit{dense score} ke rentang 0 sampai 1. Normalisasi ini diperlukan agar skor semantik memiliki skala yang sebanding dengan skor BM25.

\item \textbf{\texttt{\_minmax(scores=sparse\_scores)}}: Proses ini menormalisasi \textit{sparse score} ke rentang 0 sampai 1. Normalisasi dilakukan secara terpisah dari \textit{dense score} karena kedua skor berasal dari metode yang berbeda.

\item \textbf{\texttt{weighted\_sum(dense\_norm, sparse\_norm)}}: Proses ini menggabungkan skor \textit{dense} dan \textit{sparse} yang telah dinormalisasi. Formula yang digunakan adalah $\textit{hybrid\_score} = w_{\textit{dense}} \times \textit{dense\_norm} + w_{\textit{sparse}} \times \textit{sparse\_norm}$.

\item \textbf{\texttt{np.argsort(a=-hybrid\_scores)[:3]}}: Proses ini mengurutkan dokumen berdasarkan \textit{hybrid score} secara menurun. Tiga dokumen dengan skor tertinggi dipilih sebagai kandidat paling relevan.

\item \textbf{\texttt{\_build\_result(idx=top\_indices)}}: Proses ini mengambil metadata dokumen terpilih seperti \texttt{control\_id}, \texttt{title}, \texttt{objective}, dan \texttt{description}. Metadata beserta skor \textit{dense}, \textit{sparse}, dan \textit{hybrid} digunakan sebagai hasil akhir tahap \textit{retrieval}.
\end{enumerate}

\end{enumerate}


\subsubsection{Perancangan Komponen \textit{Rerank} pada Sistem}

Komponen \textit{reranking} pada penelitian ini menggunakan model \textit{cross-encoder} yang dipilih berdasarkan analisis perbandingan pada Subbab~2.3.4.

\vspace{1.5em}
\noindent\textbf{Pemilihan Model \textit{Reranker}}

Berdasarkan analisis perbandingan model \textit{reranker} pada Subbab~2.3.4, BAAI/bge-reranker-v2-m3 dipilih sebagai model \textit{reranker}. Model ini dipilih karena menggabungkan bukti empiris pada domain regulasi yang menunjukkan peningkatan signifikan pada kualitas peringkat, dukungan multilingual untuk lebih dari 100 bahasa yang esensial untuk \textit{query} berbahasa Indonesia, serta kemudahan implementasi sebagai model \textit{cross-encoder} yang bersifat \textit{plug-and-play}.

Perancangan komponen \textit{rerank} pada sistem dilakukan untuk meningkatkan kualitas urutan dokumen yang telah diperoleh dari tahap \textit{hybrid retrieval}. Pada rancangan ini, sistem tidak langsung menggunakan hasil \textit{hybrid retrieval} sebagai keluaran akhir, tetapi terlebih dahulu membentuk \textit{candidate pool} dan mengevaluasi ulang kandidat tersebut menggunakan model \textit{cross-encoder}. Dengan demikian, hasil akhir tidak hanya mempertimbangkan skor gabungan \textit{dense} dan \textit{sparse}, tetapi juga mempertimbangkan interaksi langsung antara \textit{query} dan dokumen kandidat.

\begin{figure}[H]
	\centering
	\includegraphics[width=\textwidth]{Gambar/3.7RerankerComponentNew.png}
	\caption{Komponen cross-encoder reranking pada sistem}
	\label{fig:reranking_komponen}
\end{figure}

\newpage

Berdasarkan alur yang ditunjukkan pada Gambar~\ref{fig:reranking_komponen}, proses \textit{rerank} dapat diuraikan ke dalam dua tahap utama, yaitu tahap \textit{hybrid retrieval} dan tahap \textit{cross-encoder reranking}.

\begin{enumerate}
\item \textbf{Tahap \textit{Hybrid Retrieval}}

Tahap \textit{hybrid retrieval} berfungsi sebagai tahap awal untuk menghasilkan kandidat dokumen yang akan dievaluasi ulang oleh \textit{reranker}. Pada tahap ini, sistem memproses \textit{query} menggunakan mekanisme \textit{dense retrieval} dan \textit{sparse retrieval}, menormalisasi skor dari kedua jalur, menggabungkannya menjadi \textit{hybrid score}, lalu memilih sejumlah kandidat teratas sebagai \textit{candidate pool}. Dalam rancangan ini, \textit{candidate pool} digunakan sebagai masukan untuk tahap \textit{cross-encoder reranking}.

\begin{enumerate}[label=\alph*., leftmargin=1.1cm]
\item \textbf{\texttt{retrieve(query=query, k=candidate\_pool)}}: Proses ini mengambil kandidat awal menggunakan komponen \textit{hybrid retrieval}. Parameter \texttt{k} menentukan jumlah kandidat yang masuk ke \textit{candidate pool} sebelum dilakukan \textit{reranking}.

\item \textbf{\texttt{\_dense\_scores(query=query)}}: Proses ini menghitung skor semantik antara \textit{query} dan seluruh dokumen menggunakan operasi \textit{dot product}. Skor yang dihasilkan merepresentasikan kemiripan makna antara kebutuhan informasi pengguna dan dokumen kontrol ISO.

\item \textbf{\texttt{score\_query(query=query)}}: Proses ini menghitung skor BM25 antara \textit{query} dan seluruh dokumen. Skor ini merepresentasikan kecocokan kata kunci berdasarkan \textit{term frequency}, IDF, dan normalisasi panjang dokumen.

\item \textbf{\texttt{\_minmax(scores=dense\_scores)}}: Proses ini menormalisasi \textit{dense score} ke rentang 0 sampai 1. Normalisasi diperlukan agar skor semantik memiliki skala yang sebanding dengan skor dari jalur \textit{sparse}.

\item \textbf{\texttt{\_minmax(scores=sparse\_scores)}}: Proses ini menormalisasi \textit{sparse score} ke rentang 0 sampai 1. Normalisasi dilakukan secara terpisah karena skor BM25 memiliki skala berbeda dari skor \textit{dense}.

\item \textbf{\texttt{weighted\_sum(dense\_norm, sparse\_norm)}}: Proses ini menggabungkan skor \textit{dense} dan \textit{sparse} yang telah dinormalisasi. Hasil penggabungan membentuk \textit{hybrid score} yang digunakan untuk menentukan kandidat awal.

\item \textbf{\texttt{np.argsort(a=-hybrid\_scores)[:k]}}: Proses ini mengurutkan dokumen berdasarkan \textit{hybrid score} secara menurun dan memilih kandidat dengan skor tertinggi. Kandidat terpilih menjadi \textit{candidate pool} untuk tahap \textit{reranking}.

\item \textbf{\texttt{hybrid\_results}}: Proses ini menyimpan daftar kandidat awal hasil \textit{hybrid retrieval}. Setiap kandidat mempertahankan metadata dokumen dan skor \textit{dense}, \textit{sparse}, serta \textit{hybrid} untuk digunakan pada tahap berikutnya.
\end{enumerate}

\item \textbf{\textit{Cross-Encoder Reranking}}

Tahap \textit{cross-encoder reranking} berfungsi untuk mengevaluasi ulang kandidat dokumen dari tahap \textit{hybrid retrieval}. Pada tahap ini, setiap kandidat diubah menjadi teks representatif, kemudian dipasangkan dengan \textit{query} dan diproses oleh model \textit{cross-encoder}. Skor \textit{cross-encoder} dinormalisasi dan dikombinasikan dengan skor \textit{hybrid} untuk menghasilkan \textit{combined score}. Skor akhir ini digunakan untuk menentukan urutan dokumen terbaik.

\begin{enumerate}[label=\alph*., leftmargin=1.1cm]
\item \textbf{\texttt{get\_model()}}: Proses ini memuat model \textit{BAAI/bge-reranker-v2-m3} dengan mekanisme \textit{lazy loading}. Model digunakan untuk menghitung relevansi pasangan \textit{query} dan dokumen secara langsung.

\item \textbf{\texttt{\_doc\_to\_rerank\_text(doc=candidate)}}: Proses ini membentuk teks kandidat yang digunakan pada tahap \textit{reranking}. Teks kandidat memuat atribut dokumen seperti judul, objektif, deskripsi, kata kunci, dan indikator audit agar konteks dokumen lebih lengkap.

\item \textbf{\texttt{make\_pairs(query, document\_texts)}}: Proses ini membentuk pasangan antara \textit{query} dan teks dokumen kandidat. Pasangan ini menjadi masukan utama bagi model \textit{cross-encoder} untuk melakukan penilaian relevansi.

\item \textbf{\texttt{model.predict(pairs=pairs)}}: Proses ini menghasilkan skor mentah \textit{cross-encoder} untuk setiap pasangan \textit{query} dan dokumen. Skor tersebut menunjukkan tingkat relevansi kandidat berdasarkan interaksi langsung antara dua teks.

\item \textbf{\texttt{\_normalize(arr=ce\_scores)}}: Proses ini menormalisasi skor \textit{cross-encoder} ke rentang 0 sampai 1. Normalisasi diperlukan agar skor \textit{cross-encoder} dapat digabungkan dengan skor \textit{hybrid}.

\item \textbf{\texttt{\_normalize(arr=hybrid\_scores)}}: Proses ini menormalisasi \textit{hybrid score} kandidat ke rentang 0 sampai 1. Skor ini digunakan sebagai sinyal relevansi awal dari tahap \textit{retrieval}.

\item \textbf{\texttt{combined\_score(alpha, ce\_norm, hybrid\_norm)}}: Proses ini menggabungkan skor \textit{cross-encoder} dan skor \textit{hybrid} yang telah dinormalisasi. Formula yang digunakan adalah $\textit{combined\_score} = \alpha \times \textit{ce\_norm} + (1 - \alpha) \times \textit{hybrid\_norm}$.

\item \textbf{\texttt{sorted(candidates, key=combined\_score)}}: Proses ini mengurutkan seluruh kandidat berdasarkan \textit{combined score} secara menurun. Kandidat dengan skor akhir tertinggi ditempatkan pada peringkat teratas.

\item \textbf{\texttt{select\_top\_n(reranked\_docs, top\_n=3)}}: Proses ini memilih tiga dokumen terbaik setelah proses \textit{reranking}. Dokumen terpilih menjadi hasil akhir yang digunakan sebagai konteks pada tahap generasi.

\item \textbf{\texttt{attach\_rerank\_metadata(doc)}}: Proses ini menambahkan informasi hasil \textit{reranking} ke metadata dokumen. Informasi tersebut mencakup \texttt{rerank\_score}, \texttt{rerank\_score\_norm}, \texttt{hybrid\_score\_norm}, dan \texttt{combined\_score} untuk mendukung transparansi hasil.
\end{enumerate}
\end{enumerate}

\section{Alur Tugas Akhir}

Penelitian ini dilaksanakan melalui serangkaian tahapan yang dirancang secara sistematis untuk menjawab hipotesis dan tujuan penelitian secara komprehensif. Setiap tahapan memiliki luaran yang saling berkaitan dan digunakan sebagai dasar pada tahapan berikutnya, sehingga keseluruhan proses penelitian dapat berjalan secara terstruktur dan berkesinambungan. Alur penelitian dimulai dari proses studi literatur dan identifikasi permasalahan, dilanjutkan dengan persiapan data dan kebutuhan penelitian, perancangan sistem dan metode yang digunakan, pelaksanaan pengujian dan eksperimen, hingga tahap analisis hasil evaluasi dan penyusunan kesimpulan akhir penelitian.

Gambar~\ref{fig:alur_penelitian} menunjukkan diagram alur penelitian secara keseluruhan yang menggambarkan hubungan antar tahapan penelitian mulai dari tahap persiapan hingga proses evaluasi akhir.

\begin{figure}[H]
	\centering
	\includegraphics[width=0.7125\textwidth,height=0.585\textheight,keepaspectratio]{Gambar/3.8AlurPenelitian.png}
	\caption{Alur penelitian}
	\label{fig:alur_penelitian}
\end{figure}

\subsection{Tinjauan Pustaka}

Langkah pertama pada penelitian ini adalah melakukan tinjauan pustaka untuk memperoleh pemahaman terkait perkembangan penelitian yang relevan dengan topik penelitian. Proses tinjauan pustaka dilakukan dengan menelusuri berbagai publikasi ilmiah, khususnya pada IEEE dan sumber ilmiah bereputasi lainnya. Beberapa kata kunci yang digunakan dalam pencarian literatur antara lain ``\textit{RAG for Auditing}'', ``\textit{Retrieval Strategy Optimization}'', ``\textit{Information Retrieval}'', dan ``\textit{Cross-Encoder Reranking}''. Kata kunci tersebut dipilih karena memiliki keterkaitan langsung dengan fokus penelitian, yaitu evaluasi strategi \textit{retrieval} dan \textit{reranking} pada \textit{pipeline} RAG.

Seperti yang telah dijelaskan pada Subbab 2.1, terdapat beberapa tema utama yang ditemukan pada penelitian-penelitian terdahulu. Variasi penelitian yang ditemukan meliputi domain penerapan sistem RAG, jenis pendekatan \textit{hybrid retrieval} yang digunakan, pemanfaatan metode \textit{cross-encoder reranking}, serta strategi pembobotan \textit{score fusion} dan penyesuaian parameter untuk memperoleh performa \textit{retrieval} yang optimal. Berdasarkan hasil tinjauan pustaka tersebut, penelitian ini dirancang dengan fokus pada evaluasi performa kombinasi \textit{hybrid retrieval} dan \textit{cross-encoder reranking} dalam meningkatkan kualitas hasil \textit{retrieval} pada dokumen ISO/IEC 27001:2022.

\subsection{Identifikasi Masalah}

Berdasarkan hasil tinjauan pustaka, telah ditemukan berbagai penelitian yang menerapkan pendekatan \textit{Retrieval-Augmented Generation} (RAG) khususnya \textit{Information Retrieval} pada domain tertentu, seperti regulasi perpajakan, dokumen keuangan, maupun kebijakan privasi. Penelitian-penelitian tersebut menunjukkan bahwa kombinasi \textit{hybrid retrieval} dan \textit{reranking} mampu meningkatkan relevansi hasil pencarian pada dokumen yang bersifat regulatif dan berbasis teks.

Namun, hingga saat penelitian ini dilakukan, masih ditemukan keterbatasan penelitian yang secara khusus membahas penerapan strategi \textit{retrieval} pada standar ISO/IEC 27001:2022, terutama pada Annex A domain A.5 hingga A.8. Karakteristik dokumen ISO 27001 memiliki tantangan tersendiri karena beberapa \textit{annex control} memiliki konteks, istilah, dan tujuan kontrol yang saling berdekatan, sehingga suatu kondisi audit dapat berkaitan dengan lebih dari satu kontrol keamanan. Hal tersebut menyebabkan proses klasifikasi dan pemetaan kontrol menjadi lebih rentan kesalahan dibandingkan domain regulasi umum lainnya.

Selain itu, pada \textit{pipeline Retrieval-Augmented Generation} (RAG), kualitas hasil \textit{retrieval} memiliki pengaruh langsung terhadap ketepatan proses \textit{generation}. Apabila dokumen atau kontrol yang dipilih pada tahap \textit{retrieval} tidak relevan, maka hasil \textit{generation} yang digunakan untuk proses pemetaan terhadap \textit{Statement of Applicability} (SoA) berpotensi menjadi kurang akurat. Oleh karena itu, penelitian ini difokuskan pada evaluasi komparatif strategi \textit{retrieval}, khususnya \textit{hybrid retrieval}, \textit{cross-encoder reranking}, serta kombinasi keduanya, untuk mengidentifikasi pendekatan yang mampu memberikan performa \textit{retrieval} paling optimal pada dokumen ISO/IEC 27001:2022. Penelitian ini tidak bertujuan untuk mengembangkan \textit{pipeline} RAG secara \textit{end-to-end}, melainkan menghasilkan rekomendasi strategi \textit{information retrieval} yang efektif untuk mendukung proses pemetaan kontrol keamanan dan implementasi sistem audit keamanan informasi berbasis ISO 27001 di masa mendatang.

\subsection{Persiapan Basis Pengetahuan Dokumen Kontrol \textit{Annex}}

Tahap persiapan basis pengetahuan dilakukan dengan mengumpulkan informasi terkait \textit{annex control} berdasarkan standar ISO/IEC 27001:2022, khususnya pada domain Annex A.5 hingga A.8. Sumber utama yang digunakan berasal dari dokumen resmi ``INTERNATIONAL STANDARD ISO/IEC 27001 Third Edition 2022-10''. Selain itu, informasi tambahan juga dikumpulkan dari berbagai sumber di internet seperti artikel implementasi keamanan informasi, dokumentasi audit, serta referensi praktik keamanan siber untuk memperkaya konteks tiap kontrol keamanan. Penambahan informasi tersebut dilakukan agar basis pengetahuan tidak hanya memuat definisi kontrol secara formal, tetapi juga mencakup istilah, konteks audit, dan variasi implementasi yang umum digunakan pada lingkungan nyata.

Seluruh informasi yang telah dikumpulkan kemudian direpresentasikan dalam format JSON (\textit{JavaScript Object Notation}), di mana setiap \textit{annex control} disimpan sebagai satu entitas data terstruktur. Struktur JSON mencakup atribut seperti identitas kontrol, deskripsi, panduan implementasi, kata kunci, indikator audit, aset terkait, dan prinsip keamanan dalam Bahasa Indonesia maupun Bahasa Inggris. Proses penyusunan dan pengisian atribut dilakukan secara manual oleh penulis untuk menjaga kesesuaian konteks dengan standar ISO 27001, serta dibantu dengan pemanfaatan \textit{Generative AI} untuk menghasilkan variasi istilah, pengayaan kata kunci, dan penyusunan atribut pendukung yang lebih lengkap.

\subsection{Persiapan \textit{Dataset} Variasi \textit{Query}}

Tahap persiapan \textit{dataset query} dilakukan dengan menyusun \textit{query} yang dibentuk dari dokumen asli milik Direktorat Teknologi Informasi Universitas Gadjah Mada (DTI UGM). Proses penyusunan \textit{query} dilakukan melalui pembacaan dan penelaahan manual terhadap dokumen kebijakan, prosedur, standar operasional, tata kelola, serta regulasi internal yang berkaitan dengan keamanan informasi dan pengelolaan teknologi informasi. Dari hasil pembacaan tersebut, kemudian diidentifikasi potongan informasi, kebutuhan audit, maupun pernyataan kebijakan yang berpotensi merepresentasikan kebutuhan pencarian informasi pada proses audit ISO/IEC 27001:2022. Berdasarkan konteks tersebut, \textit{query} kemudian dirumuskan secara manual dengan mempertahankan istilah, struktur bahasa, dan gaya penyampaian yang mendekati kondisi nyata pada lingkungan organisasi DTI UGM.

Setelah \textit{query} berhasil dibentuk, dilakukan proses penelaahan untuk menentukan \textit{ground truth annex control} yang sesuai terhadap masing-masing \textit{query}. Penentuan \textit{ground truth} dilakukan dengan mencocokkan konteks \textit{query} terhadap kontrol ISO/IEC 27001:2022 Annex A.5 hingga A.8 berdasarkan relevansi isi, tujuan pengendalian, serta cakupan keamanan informasi yang dibahas dalam dokumen sumber. Setiap \textit{query} dapat memiliki keterkaitan dengan satu atau lebih kontrol \textit{annex} sesuai kompleksitas konteks yang ditemukan.

\subsection{Perancangan \textit{Hybrid Retrieval} dan Pemilihan Parameter Optimal}

Perancangan metode \textit{hybrid retrieval} pada penelitian ini dilakukan dengan mengacu pada beberapa penelitian terdahulu, khususnya penelitian oleh Rustamov et al. dan Liu et al. yang menerapkan kombinasi \textit{sparse retrieval} menggunakan BM25 dan \textit{dense retrieval} berbasis model MiniLM. Pada penelitian ini, BM25 digunakan untuk menangkap kesesuaian istilah secara eksplisit, sedangkan \textit{dense retrieval} menggunakan model paraphrase-multilingual-MiniLM-L12-v2 untuk memahami keterkaitan makna semantik antara \textit{query} dan dokumen \textit{annex control}. Penggunaan model ini merupakan pengembangan dari penelitian sebelumnya yang banyak menggunakan model all-MiniLM-L6-v2 sebagai \textit{baseline dense retrieval}.

Pada penelitian ini dipilih model \textit{multilingual} karena kumpulan \textit{query} dan dokumen berasal dari lingkungan institusi pendidikan di Indonesia yang didominasi penggunaan Bahasa Indonesia serta istilah teknis berbahasa Inggris secara bersamaan. Selain itu, arsitektur L12 dengan 12 \textit{layer} dipilih karena memiliki kapasitas representasi semantik yang lebih baik dibandingkan varian L6, sehingga diharapkan lebih mampu menangkap hubungan makna pada \textit{query} audit yang bersifat implisit, panjang, dan kontekstual. Dengan demikian, model ini dinilai lebih sesuai untuk skenario \textit{retrieval} dokumen keamanan informasi pada domain DTI UGM yang memiliki karakteristik bahasa campuran dan konteks kebijakan yang kompleks.

Setelah metode \textit{hybrid retrieval} dirancang, tahap selanjutnya adalah melakukan pencarian parameter bobot \textit{score fusion} yang optimal melalui penyesuaian nilai $\alpha$ (\textit{alpha scoring}). Parameter $\alpha$ digunakan untuk mengatur proporsi kontribusi skor antara \textit{sparse retrieval} dan \textit{dense retrieval} pada proses penggabungan skor akhir. Pengujian dilakukan terhadap seluruh variasi \textit{query} yang digunakan dalam penelitian untuk menentukan kombinasi bobot yang mampu memberikan performa \textit{retrieval} terbaik sebelum integrasi metode \textit{reranking}. Hasil pencarian parameter ini kemudian digunakan sebagai konfigurasi dasar pada tahap eksperimen berikutnya agar proses \textit{score fusion} antara \textit{hybrid retrieval} dan \textit{cross-encoder reranking} dapat dilakukan secara konsisten dan terukur.

\subsection{Perancangan \textit{Reranking} dan Pemilihan Parameter Optimal}

Tahap perancangan \textit{reranking} dilakukan sebagai pengembangan dari rancangan \textit{hybrid retrieval} yang telah dibangun pada tahap sebelumnya. Pada proses ini, hasil \textit{retrieval} awal dari metode \textit{hybrid retrieval} digunakan sebagai \textit{candidate pool} yang akan diproses kembali untuk meningkatkan kualitas peringkat dokumen. Dengan pendekatan ini, sistem tidak melakukan proses \textit{reranking} terhadap seluruh dokumen pada korpus, melainkan hanya pada sejumlah dokumen dengan skor \textit{retrieval} tertinggi hasil dari proses \textit{hybrid retrieval}. Pendekatan tersebut dipilih untuk menjaga efisiensi komputasi sekaligus mempertahankan kualitas relevansi hasil pencarian.

Metode \textit{reranking} yang digunakan pada penelitian ini adalah pendekatan \textit{cross-encoder} menggunakan model BAAI/bge-reranker-v2-m3 yang mendukung pemrosesan multibahasa, termasuk Bahasa Indonesia. Pemilihan model BGE \textit{reranker} juga didasarkan pada beberapa penelitian sebelumnya oleh Rustamov et al. menggunakan seri model BGE sebagai pendekatan \textit{retrieval} dan \textit{reranking} karena memiliki performa \textit{semantic matching} yang baik pada tugas \textit{information retrieval} modern. Selain itu, penelitian oleh Elkiran-Rasheed menunjukkan bahwa model BAAI/bge-reranker-base mampu menghasilkan nilai akurasi terbaik dibandingkan beberapa model \textit{reranker} lainnya pada proses evaluasi \textit{retrieval}.

Berdasarkan pertimbangan tersebut, penelitian ini menggunakan varian BAAI/bge-reranker-v2-m3 karena memiliki kemampuan \textit{multilingual} dan kapasitas pemahaman konteks yang lebih sesuai dengan karakteristik dokumen kebijakan serta \textit{query} audit pada lingkungan institusi pendidikan di Indonesia. Setelah rancangan \textit{reranking} dibangun, dilakukan pengujian terhadap beberapa variasi jumlah \textit{candidate pool} yang diambil dari hasil \textit{hybrid retrieval} untuk menentukan nilai top-$n$ yang paling optimal untuk proses \textit{reranking} dilakukan. Pengujian ini bertujuan memperoleh keseimbangan antara kualitas relevansi dan efisiensi komputasi, sehingga proses \textit{score fusion} antara \textit{hybrid retrieval} dan \textit{cross-encoder reranking} dapat menghasilkan performa \textit{retrieval} yang optimal.

\subsection{Pengujian Kombinasi dan Pemilihan \textit{Scoring} Optimal}

Setelah perancangan metode \textit{hybrid retrieval} dan \textit{cross-encoder reranking} selesai dilakukan, tahap selanjutnya adalah melakukan pengujian kombinasi kedua metode melalui mekanisme pembobotan \textit{score fusion}. Pada tahap ini, pengujian dilakukan secara komprehensif dengan memvariasikan parameter $\alpha$ (\textit{alpha scoring}) untuk menentukan proporsi skor yang paling optimal antara hasil \textit{hybrid retrieval} dan skor dari \textit{reranking}, dengan urutan pembobotan \textit{hybrid}:\textit{cross-encoder}. Dalam kerangka ini, pengujian \textit{hybrid retrieval} secara independen direpresentasikan oleh parameter pembobotan 1:0, pengujian \textit{reranker} secara independen direpresentasikan oleh parameter pembobotan 0:1, sedangkan konfigurasi dengan pembobotan lainnya merepresentasikan kombinasi kedua metode.

Pada tahap ini, \textit{hybrid retrieval} menggunakan parameter bobot $\alpha$ paling optimal antara komponen \textit{dense retrieval} dan \textit{sparse retrieval} yang diperoleh dari tahap sebelumnya, sedangkan \textit{reranking} menggunakan jumlah \textit{candidate pool} terbaik. Proses eksperimen dilakukan pada seluruh \textit{query} yang dibentuk dari dokumen kebijakan Direktorat Teknologi Informasi (DTI) UGM. Setiap kombinasi parameter dievaluasi menggunakan metrik \textit{Information Retrieval} seperti Hit@3, Recall@3, MAP@3, dan NDCG@3 untuk mengidentifikasi konfigurasi bobot \textit{score fusion} yang memberikan performa paling optimal. Pengujian ini bertujuan untuk mengevaluasi kualitas performa masing-masing pendekatan baik secara independen maupun dalam bentuk kombinasi, sehingga dapat diketahui kontribusi dan karakteristik tiap metode terhadap hasil \textit{retrieval} sekaligus ditentukan pendekatan kombinasi yang paling efektif dalam meningkatkan relevansi hasil pencarian pada proses pemetaan \textit{annex control} ISO/IEC 27001:2022.

\subsection{Analisis Evaluasi Hasil}

Tahap analisis evaluasi hasil dilakukan setelah seluruh proses pengujian terhadap metode \textit{hybrid retrieval}, \textit{cross-encoder reranking}, dan kombinasi bobot \textit{score fusion} selesai dilaksanakan. Pada tahap ini, hasil evaluasi dibandingkan untuk menganalisis performa masing-masing pendekatan, baik secara independen maupun dalam bentuk kombinasi. Analisis dilakukan menggunakan beberapa metrik \textit{Information Retrieval} seperti Hit@3, Recall@3, MAP@3, dan NDCG@3 untuk melihat tingkat relevansi dan kualitas peringkat dokumen yang dihasilkan oleh setiap metode. Selain itu, tahap ini juga digunakan untuk mengamati pengaruh perubahan parameter terhadap performa \textit{retrieval}, seperti nilai $\alpha$ pada \textit{hybrid retrieval}, ukuran \textit{candidate pool}, serta proporsi \textit{score fusion} antara \textit{retrieval} dan \textit{reranking}.

Melalui analisis tersebut, diharapkan dapat terlihat pola performa akurasi dari tiap metode terhadap \textit{query} yang dibentuk dari dokumen kebijakan DTI UGM. Hasil analisis ini diharapkan dapat memberikan rekomendasi konfigurasi \textit{retrieval} yang paling sesuai untuk implementasi sistem \textit{Information Retrieval} pada domain audit keamanan informasi berbasis ISO/IEC 27001:2022.

\subsection{Penulisan Laporan}

Setelah seluruh tahapan penelitian selesai dilaksanakan, dilakukan proses penyusunan laporan sebagai bentuk dokumentasi akhir penelitian. Laporan ini memuat uraian lengkap mengenai tahapan penelitian yang telah dilakukan, implementasi metode yang digunakan, hasil eksperimen dan pengujian yang diperoleh, proses evaluasi dan interpretasi hasil, serta rangkuman temuan penelitian.






